\section{Problem formulation}\label{sec:formulation}

\paragraph{Video and chunks.}
Let $V$ be a video of duration $T$ seconds. It is partitioned into $n$
consecutive chunks $c_i=[u_i,v_i)$, $i=1,\dots,n$, with $u_1=0$, $v_n=T$ and
$v_i=u_{i+1}$ (the construction is given in \cref{sec:chunking}).

\paragraph{Perception.}
A fixed perception function $\phi$ maps each chunk, through its sampled frames
and the transcript segments attached to it, to an observation
$o_i=\phi(c_i)$. The function $\phi$ is \emph{memory-agnostic}: $o_i$ depends
on $c_i$ only, not on any memory state. The observations
$o_{1:n}$ are computed once and frozen.

\paragraph{Memory.}
A memory organisation is a pair $(U,R)$ of an update rule and a retrieval rule.
Starting from $M_0=\varnothing$, the update rule produces
\begin{equation}
  M_i = U(M_{i-1}, o_i), \qquad i=1,\dots,n,
\end{equation}
subject to the budget constraint
\begin{equation}\label{eq:budget}
  \toklen{M_i} \le B \quad\text{for all } i,
  \qquad \toklen{M} := \bigl|\tau\bigl(\sigma(M)\bigr)\bigr|,
\end{equation}
where $\sigma$ serialises the memory to text and $\tau$ is the tokenizer of
the answering model. Measuring the budget in answerer tokens makes it the
quantity that actually constrains answering, and makes budgets comparable
across organisations with different record formats.

\paragraph{Provenance.}
Every record $r\in M$ carries a provenance set $\pi(r)\subseteq\{1,\dots,n\}$
of the chunks it was derived from. The update rule must preserve provenance:
a record created by merging or compressing records $r_1,\dots,r_k$ receives
$\pi=\bigcup_j\pi(r_j)$. The time support of a set of records $C$ is
\begin{equation}
  I(C) = \bigcup_{r\in C}\;\bigcup_{i\in\pi(r)} [u_i,v_i).
\end{equation}

\paragraph{Answering.}
For a question $q$, the retrieval rule selects $R(q,M_n)\subseteq M_n$ with
$\toklen{R(q,M_n)}\le B_R$ (the one exception is C1, whose memory is a single
summary of size at most $B$ that is always given whole). A fixed answering model
$A$ returns an answer and the
subset of retrieved records it cites: $(\hat a,\,C)=A\bigl(q,R(q,M_n)\bigr)$,
$C\subseteq R(q,M_n)$. The answerer has no access to the video.

\paragraph{Questions.}
Each question $q$ in the evaluation set $\mathcal Q$ has a reference answer
$a^\ast(q)$, a type $\kappa(q)$ (past event, object change, cross-segment
relation, attribute recall, speech-grounded), ground-truth evidence intervals
$G(q)$, and an evidence time $e(q)$ (the latest moment needed to answer). Its
\emph{evidence age} is $\age(q)=T-e(q)$: the time elapsed between the evidence
and the end of the video, i.e.\ how long the memory had to retain it.

\paragraph{Evaluation quantities.}
For organisation $m$ at budget $B$:
\begin{align}
  y_m(q) &= \ind\bigl[\hat a \text{ is correct for } a^\ast(q)\bigr], &
  \Acc_m(B) &= \tfrac{1}{|\mathcal Q|}\textstyle\sum_{q} y_m(q),\\
  g_m(q) &= \ind\bigl[I(C)\cap G(q)\neq\varnothing\bigr], &
  \mathrm{tIoU}_m(q) &= \frac{|I(C)\cap G(q)|}{|I(C)\cup G(q)|},\\
  p_m(q) &= \frac{|I(C)\cap G(q)|}{|I(C)|}, &&
\end{align}
where $|\cdot|$ of a union of intervals is its total length and $p_m(q)=0$ when
nothing is cited. Grounding $g$ alone rewards citing records with long support
(an entity seen throughout the video overlaps every interval); the
\emph{grounding precision} $p$ and tIoU penalise it. Retrieval recall is $g$
computed with $R(q,M_n)$ in place of $C$. \emph{Evidence coverage}, used as a
proxy for fact survival, is $g$ computed with the non-entity records of $M_n$ in
place of $C$: some event, state change, relation or note derived from the
evidence chunks is still stored. It is a proxy because it shows that material
from those chunks survives, not that the specific fact does; the failure
attribution (\cref{sec:attribution}) checks the latter by hand.

\paragraph{Controlled comparison.}
The organisations compared later (\cref{tab:conditions}) differ only in
$(U,R)$. The quantities $\phi$, $o_{1:n}$, $\tau$, $B$, $B_R$ and $A$ are
shared. This is the formal content of the requirement that comparisons use
``the same observations and a similar memory budget''.

\begin{table}[ht]
\centering\small
\caption{Memory organisations to be compared (Stage~5). All share the
observations, the answerer and the retrieval allowance $B_R$.}
\label{tab:conditions}
\begin{tabularx}{\linewidth}{@{}llX@{}}
\toprule
ID & Organisation & Description \\
\midrule
C1 & Rolling summary & One text record, rewritten after each chunk under
  budget $B$; the whole summary is given to the answerer. \\
C2 & Flat log (FIFO) & One text record per observation; oldest evicted when
  over budget; embedding retrieval. \\
C3 & Structured memory & Typed records (entities, events, state changes,
  relations), entity resolution, compression, salience-based deletion;
  hybrid retrieval. \\
C4 & Unbounded flat log & C2 without eviction: an upper bound for the given
  perception. \\
\bottomrule
\end{tabularx}
\end{table}

\paragraph{A consequence for grounding.}
For C1 the memory is a single record whose provenance is every chunk seen so
far, so $I(C)$ covers most of the video and $g_{\mathrm{C1}}\approx 1$
trivially, and coverage is equally uninformative. Grounding, coverage and failure
attribution are therefore reported only for C2--C4; for C1 answer accuracy is
reported, overall, per question type and against evidence age.
